\subsection{比较关系：I. 弱关系}

在下文中，我们用 V 表示赋值域 K 上的一个赋范向量空间，并用 \(\mathcal{H}(\mathfrak{F}, V)\) 表示取值于 V 且各自定义在 E 的属于滤子基 F 的某个子集上的函数所成的集。我们将要在这些函数之间定义的关系，相对于以 F 为基的滤子具有局部性：我们来明确这句话的含义。设 f 与 g 是 \(\mathcal{H}(\mathfrak{F}, V)\) 中的两个函数，回忆一下，关系“存在一个集 \(Z \in F\) ，使 f 与 g 在 Z 上都有定义且相等”是一个等价关系 \(R_{\infty}\) ，定义在 \(\mathcal{H}(\mathfrak{F}, V)\) 上 (Gen.~Top., I, p.~66)。在此情形下，若含 \(\mathcal{H}(\mathfrak{F}, V)\) 中一个函数 f 的关系 S（关于 f）与等价关系 \(R_{\infty}\) 相容 (Set Theory, II, p.~117)，我们就说它（沿 F）关于 f 具有局部性；我们知道，如果 \(\tilde{f}\) 是 f 沿 F 的芽，即 f 模 \(R_{\infty}\) 的等价类（商集 \(\mathcal{H}_{\infty}(\mathfrak{F}, V) = \mathcal{H}(\mathfrak{F}, V)/R_{\infty}\) 中的一个元素），那么由 S 过渡到商，可以导出一个 \(\tilde{f}\) 与 S 的其余变元之间的关系；反之，每个这种形式的关系也定义一个关于 f 具有局部性的关系。

例. 设 f 与 g 是 \(\mathcal{H}(\mathfrak{F},\mathbf{R})\) 中的两个函数，关系“存在一个 \(X \in F\) ，使 f 与 g 都在 X 上有定义，且 \(f(t) \leqslant g(t)\) 对每个 \(t \in X\) 成立”关于 f 与 g 具有局部性。我们把（对 f 与 g）过渡到商所得到的关系记作 \(\tilde{f} \leqslant \tilde{g}\) ；我们指出，若 \(\tilde{f} \leqslant \tilde{g}\) ，则存在一个函数 \(f_{1} \in \tilde{f}\) 与一个函数 \(g_{1} \in \tilde{g}\) ，都定义在整个 E 上，使得 \(f_{1}(t) \leqslant g_{1}(t)\) 对一切 \(t \in E\) 成立。

注. 1) 设 \(V_{i}\) ( \(1 \leqslant i \leqslant n\) ) 是 K 上的 n 个赋范向量空间，\(\varphi\) 是定义在 \(V_{1} \times V_{2} \times \cdots \times V_{n}\) 上、取值于 V 的函数；沿 \(R_{\infty}\) 过渡到商，函数 \(\varphi\) 便定义了一个从

\[
\mathcal {H} _ {\infty} (\mathfrak {F}, V _ {1}) \times \dots \times \mathcal {H} _ {\infty} (\mathfrak {F}, V _ {n})
\]

到 \(\mathcal{H}_{\infty}(\mathfrak{F}, V)\) 中的映射，通常记作 \(\varphi(\tilde{\mathbf{f}}_1, \ldots, \tilde{\mathbf{f}}_n)\) (Gen.~Top., I, p.~67)。例如，取 \(\varphi\) 为映射 \((\mathbf{x}, \mathbf{y}) \mapsto \mathbf{x} + \mathbf{y}\) 与 \(\mathbf{x} \mapsto \mathbf{x}\lambda\) ( \(\lambda \in \mathbf{K}\) )，便对任意两个芽 \(\tilde{\mathbf{f}}\) 、 \(\tilde{\mathbf{g}}\) （均属于 \(\mathcal{H}_{\infty}(\mathfrak{F}, V)\) ）定义了元素 \(\tilde{\mathbf{f}} + \tilde{\mathbf{g}}\) 与 \(\tilde{\mathbf{f}}\lambda\) ，并且立即可验证，合成法则 \((\tilde{\mathbf{f}}, \tilde{\mathbf{g}}) \mapsto \tilde{\mathbf{f}} + \tilde{\mathbf{g}}\) 与 \((\lambda, \tilde{\mathbf{f}}) \mapsto \tilde{\mathbf{f}}\lambda\) 在 \(\mathcal{H}_{\infty}(\mathfrak{F}, V)\) 上定义了域 \(\mathbf{K}\) 上的一个向量空间结构；在这个空间中，\(\tilde{0}\) 是由等于 \(0\) （于 \(\mathfrak{F}\) 的某个集上）的函数组成的类，而 \(-\tilde{\mathbf{f}}\) 是由等于 \(-\mathbf{f}\) （于 \(\mathfrak{F}\) 的某个集上）的函数组成的类。同样，若 \(V\) 是 \(K\) 上的一个代数，则对 \(\mathcal{H}_{\infty}(\mathfrak{F}, V)\) 定义第二个内合成法则 \((\tilde{\mathbf{f}}, \tilde{\mathbf{g}}) \mapsto \tilde{\mathbf{f}}\tilde{\mathbf{g}}\) （取 \(\varphi(\mathbf{x}, \mathbf{y}) = \mathbf{x}\mathbf{y}\) ）；它与前两个法则一起，定义了 \(K\) 上的一个代数结构于 \(\mathcal{H}_{\infty}(\mathfrak{F}, V)\) 上；若 \(V\) 具有单位元 \(\mathbf{e}\) ，则 \(\mathcal{H}_{\infty}(\mathfrak{F}, V)\) 以类 \(\tilde{\mathbf{e}}\) 为单位元，该类由等于 \(\mathbf{e}\) 于 \(\mathfrak{F}\) 的某个集上的函数组成；为使 \(\tilde{\mathbf{f}}\) 在 \(\mathcal{H}_{\infty}(\mathfrak{F}, V)\) 中可逆，必须且只需：对某个 \(\mathbf{f} \in \tilde{\mathbf{f}}\) ，存在一个 \(Z \in \mathfrak{F}\) ，使得 \(\mathbf{f}(t)\) 在 \(V\) 中对每个 \(t \in Z\) 都可逆（这时，类 \(\tilde{\mathbf{f}}\) 中每个函数都满足这一条件）。

\begin{enumerate}
\def\labelenumi{\arabic{enumi})}
\setcounter{enumi}{1}
\tightlist
\item
  沿用同样的记号，设 \(\psi\) 是 \(\prod_{i=1}^{n} V_i\) 的一个子集到 \(V\) 中的映射；我们以 \(\psi(\mathbf{f}_1, \mathbf{f}_2, \ldots, \mathbf{f}_n)\) 表示这样的函数，它等于 \(\psi(\mathbf{f}_1(t), \mathbf{f}_2(t), \ldots, \mathbf{f}_n(t))\) 于每个点 \(t \in E\) （在该点诸 \(\mathbf{f}_i(t)\) 有定义，且点 \((\mathbf{f}_i(t))\) 属于 \(\psi\) 的定义集）\(^1\) 。例如，\(\mathbf{f} + \mathbf{g}\) 就是等于 \(\mathbf{f}(t) + \mathbf{g}(t)\) 于每个点 \(t \in E\) （\(\mathbf{f}\) 与 \(\mathbf{g}\) 在该点都有定义）的函数。注意，映射 \((\mathbf{f}, \mathbf{g}) \mapsto \mathbf{f} + \mathbf{g}\) 不是 \(\mathcal{H}(\mathfrak{F}, V)\) 上的一个群律，因为若 \(\mathbf{f}\) 不在整个 \(E\) 上有定义，则不存在函数 \(\mathbf{g} \in \mathcal{H}(\mathfrak{F}, V)\) 使 \(\mathbf{f} + \mathbf{g} = 0\) 。
\end{enumerate}

定义 1. 给定两个实函数 \(f, g\) ，它们属于 \(\mathcal{H}(\mathfrak{F}, \mathbf{R})\) 且 \(\geqslant 0\) 于 \(\mathfrak{F}\) 的某个集上，我们说 \(f\) 受 \(g\) 控制，或者说 \(g\) 控制 \(f\) （沿 \(\mathfrak{F}\) ），并记作 \(f \preccurlyeq g\) 或 \(g \succcurlyeq f\) ，如果存在一个集 \(X \in \mathfrak{F}\) 和一个数 \(k > 0\) ，使得 \(f(t) \leqslant k g(t)\) 对每个 \(t \in X\) 成立（换句话说，如果存在 \(k > 0\) 使 \(\tilde{f} \leqslant k \tilde{g}\) ）

给定两个赋范向量空间 \(V_{1}, V_{2}\) ，以及两个函数 \(f_{1}, f_{2}\) ，分别属于 \(\mathcal{H}(\mathfrak{F}, V_{1})\) 与 \(\mathcal{H}(\mathfrak{F}, V_{2})\) ，我们说 \(f_{1}\) 受 \(f_{2}\) 控制（沿 \(F\) ），并记作 \(f_{1} \preccurlyeq f_{2}\) 或 \(f_{2} \succcurlyeq f_{1}\) ，如果 \(\|f_{1}\| \preccurlyeq \|f_{2}\|\) 。

关系 \(f_{1} \preccurlyeq f_{2}\) 显然关于 \(f_{1}\) 与 \(f_{2}\) 具有局部性；因此它等价于过渡到商所得到的关系 \(\tilde{f}_{1} \preccurlyeq \tilde{f}_{2}\) 。当 f 与 g 是实函数时，应当注意不要混淆关系 \(\tilde{f} \preccurlyeq \tilde{g}\) 与 \(\tilde{f} \leqslant \tilde{g}\) 。

注意，对每个标量 \(\lambda \neq 0\) ，关系 \(\mathbf{f}_1 \preccurlyeq \mathbf{f}_2\lambda\) 等价于 \(\mathbf{f}_1 \preccurlyeq \mathbf{f}_2\) 。若 \(\mathbf{f}_1 \preccurlyeq \mathbf{f}_2\) ，则存在一个集 \(X \in \mathfrak{F}\) ，使得 \(\mathbf{f}_1(x) = 0\) 对每个点 \(x \in X\) （在该点 \(\mathbf{f}_2(x) = 0\) ）成立。

例. 1) 关系 \(\mathbf{f} \preccurlyeq 1\) 表示 \(\mathbf{f}\) 在 \(\mathfrak{F}\) 的某个集上有界。

\begin{enumerate}
\def\labelenumi{\arabic{enumi})}
\setcounter{enumi}{1}
\item
  对每个函数 \(\mathbf{f}\) （属于 \(\mathcal{H}(\mathfrak{F}, V)\) ）以及每个标量 \(\lambda \neq 0\) ，有 \(\mathbf{f} \preccurlyeq \mathbf{f}\lambda\) 。
\item
  当 \(x\) 趋于 \(+\infty\) 时，有 \(\sin^2 x \preccurlyeq \sin x\) 。
\item
  当 \((x,y)\) 趋于 \((0,0)\) （在 \(\mathbf{R}^2\) 中）时，有
\end{enumerate}

\[
x y \preccurlyeq x ^ {2} + y ^ {2}.
\]

下列命题是定义 1 的直接推论：

命题 1. 若 \(f, g, h\) 是 \(\mathcal{H}(\mathfrak{F}, \mathbf{R})\) 中的三个函数，则关系 \(f \preccurlyeq g\) 与 \(g \preccurlyeq h\) 蕴含 \(f \preccurlyeq h\) 。

命题 2. 设 \(\mathbf{f}_1, \mathbf{f}_2\) 是 \(\mathcal{H}(\mathfrak{F}, V)\) 中的两个函数，\(g\) 是 \(\mathcal{H}(\mathfrak{F}, \mathbf{R})\) 中的一个函数。关系 \(\mathbf{f}_1 \preccurlyeq g\) 与 \(\mathbf{f}_2 \preccurlyeq g\) 蕴含 \(\mathbf{f}_1 + \mathbf{f}_2 \preccurlyeq g\) 。

进一步：

命题 3. 设 \(V_{1}, V_{2}, V\) 是同一赋值域上的三个赋范空间，\((\mathbf{x}, \mathbf{y}) \mapsto [\mathbf{x}, \mathbf{y}]\) 是从 \(V_{1} \times V_{2}\) 到 V 中的一个双线性映射。若 \(f_{1}\) 与 \(f_{2}\) 分别是 \(\mathcal{H}(\mathfrak{F}, V_{1})\) 与 \(\mathcal{H}(\mathfrak{F}, V_{2})\) 中的函数，且 \(g_{1}, g_{2}\) 是 \(\mathcal{H}(\mathfrak{F}, \mathbf{R})\) 中使 \(f_{1} \preccurlyeq g_{1}\) 与 \(f_{2} \preccurlyeq g_{2}\) 成立的两个函数，则 \([f_{1}.f_{2}] \preccurlyeq g_{1}g_{2}\) 。

事实上，(Gen.~Top., IX, p.~173, th. 1) 存在一个数 \(a > 0\) ，使得 \(\|[\mathbf{f}_1.\mathbf{f}_2]\| \leqslant a\|\mathbf{f}_1\|\|\mathbf{f}_2\|\) 。

推论. 若 V 是一个赋范代数，\(f_{1}, f_{2}\) 是 \(\mathcal{H}(\mathfrak{F}, V)\) 中的两个函数，且 \(g_{1}, g_{2}\) 是 \(\mathcal{H}(\mathfrak{F}, \mathbf{R})\) 中的两个函数，则关系 \(f_{1} \preccurlyeq g_{1}\) 、 \(f_{2} \preccurlyeq g_{2}\) 蕴含 \(f_{1}f_{2} \preccurlyeq g_{1}g_{2}\) 。

由命题 1，关系 \(f \preccurlyeq g\) （\(\mathcal{H}(\mathfrak{F}, V)\) 中函数之间的）是传递的；又因它是自反的，关系“f \(\preccurlyeq g\) 且 \(g \preccurlyeq f\) ”是 \(\mathcal{H}(\mathfrak{F}, V)\) 上的一个等价关系 (Set Theory, II, p.~113)。

定义 2. 给定两个函数 \(\mathbf{f}\) 、 \(\mathbf{g}\) （属于 \(\mathcal{H}(\mathfrak{F}, V)\) ），我们说 \(\mathbf{f}\) 与 \(\mathbf{g}\) 相似（沿 \(\mathfrak{F}\) ），并记作 \(\mathbf{f} \asymp \mathbf{g}\) ，如果 \(\mathbf{f} \preccurlyeq \mathbf{g}\) 且 \(\mathbf{g} \preccurlyeq \mathbf{f}\) 。

对每个标量 \(\lambda \neq 0\) ，关系 \(\mathbf{f} \asymp \mathbf{g}\) 等价于 \(\mathbf{f} \asymp \mathbf{g}\lambda\) 。它蕴含存在一个集 \(X \in \mathfrak{F}\) ，使得 \(X\) 中使 \(\mathbf{f}(x) = 0\) 的点所成的子集与 \(X\) 中使 \(\mathbf{g}(x) = 0\) 的点所成的子集相同。

例. 1) 对实函数 \(f \in \mathcal{H}(\mathfrak{F}, \mathbf{R})\) 而言，关系 \(f \asymp 1\) 表示存在两个数 \(a > 0\) 、 \(b > 0\) ，使得 \(a \leqslant |f(x)| \leqslant b\) 于 \(\mathfrak{F}\) 的某个集上成立，或者表示函数 \(\log |f|\) 在 \(\mathfrak{F}\) 的某个集上有界：这时称 \(f\) 在 \(\mathfrak{F}\) 的某个集上对数有界。

\begin{enumerate}
\def\labelenumi{\arabic{enumi})}
\setcounter{enumi}{1}
\item
  设 V 是非离散赋值域 K 上的一个赋范空间，\(\mathbf{f}(x)=\mathbf{a}_{0}x^{n}+\mathbf{a}_{1}x^{n-1}+\cdots+\mathbf{a}_{n}\) 是变元 \(x\in K\) 的、以 V 为系数集的多项式，且 \(a_{0}\neq0\) 。对每个向量 \(b\neq0\) ，有 \(\mathbf{f}(x)\asymp\mathbf{b}x^{n}\) （当 \(|x|\) 趋于 \(+\infty\) ）。
\item
  我们已经看到 \(\sin^2 x \preccurlyeq \sin x\) （当 \(x\) 趋于 \(+\infty\) 时），但 \(\sin^2 x \asymp \sin x\) 并不成立，尽管这两个函数在相同的点处为零。
\item
  有 \(x^{2} + xy + y^{2} \asymp x^{2} + y^{2}\) （当 \((x, y)\) 趋于 \((0, 0)\) ，在 \(\mathbf{R}^2\) 中），但 \(xy \asymp x^{2} + y^{2}\) 不成立。
\end{enumerate}

由 V, p.~213 的命题 3 立即得出：若 \(f_1\) 、 \(f_2\) 、 \(g_1\) 、 \(g_2\) 是 \(\mathcal{H}(\mathfrak{F},\mathrm{K})\) 中的函数（K 为任一赋值域），则关系 \(f_1 \asymp g_1\) 与 \(f_2 \asymp g_2\) 蕴含 \(f_1f_2 \asymp g_1g_2\) 。

我们指出，与此相反，关系 \(f_{1} \asymp g_{1}\) 与 \(f_{2} \asymp g_{2}\) 并不蕴含 \(f_{1} + f_{2} \asymp g_{1} + g_{2}\) ，例如 \(f_{1}(x) = g_{1}(x) = x^{2}\) 、 \(f_{2}(x) = -(x^{2} + x)\) 、 \(g_{2}(x) = -(x^{2} - 1)\) （当实变量 \(x\) 趋于 \(+\infty\) ）便说明了这一点。

比较关系 \(f \preccurlyeq g\) 、 \(f \asymp g\) 称为弱关系。我们说 \(\mathcal{H}(\mathfrak{F}, V)\) 中的两个函数 f 、 g 是弱可比较的，如果它们满足关系 \(f \preccurlyeq g\) 、 \(g \preccurlyeq f\) 中的至少一个。

注. 1) \(\mathcal{H}(\mathfrak{F},\mathbf{R})\) 中的两个函数未必弱可比较，例如函数 1 与 \(x\sin x\) （当 \(x\) 趋于 \(+\infty\) ）便说明了这一点。

\begin{enumerate}
\def\labelenumi{\arabic{enumi})}
\setcounter{enumi}{1}
\tightlist
\item
  用 \(\mathbf{R}_0\) 表示关系 \(\mathbf{f} \asymp \mathbf{g}\) （在 \(\mathcal{H}(\mathfrak{F}, V)\) 上），用 \(\mathcal{H}_0(\mathfrak{F}, V)\) 表示商集 \(\mathcal{H}(\mathfrak{F}, V)/\mathrm{R}_0\) ；注意关系 \(\mathrm{R}_{\infty}\) 蕴含 \(\mathbf{R}_0\) 。对关系 \(\mathbf{f} \preccurlyeq \mathbf{g}\) 过渡到商，由 V, p.~213 的命题 1 便得到 \(\mathcal{H}_0(\mathfrak{F}, V)\) 上的一个序关系 (Set Theory, III, p.~134)；前一例子表明 \(\mathcal{H}_0(\mathfrak{F}, V)\) 并不被这个关系全序化。
\end{enumerate}

